\begin{textsample}{POS Dim 6 – mg – Score 10.00 – mg/mg\_em\_94.txt}  \label{ex:f6_pos_243}
4.9.3 Área de Ciências da Natureza e suas Tecnologias O Itinerário proposto para a área de Ciência da Natureza está alinhado com as orientações pedagógicas propostas na Formação Geral Básica.
Essas devem estar fundamentadas nas concepções educativas que contemplem o pluralismo de ideias, e, sobretudo, sejam construídas coletivamente, pautadas no contexto ao qual a escola está inserida.
Sugere -se o uso das abordagens CTSA enquanto ferramenta metodológica para integrar os eixos \textbf{estruturantes} dos Itinerários e nessa mesma perspectiva o Ensino de Ciências por Investigação associado às sequências didáticas investigativas em que o estudante deixa de mero expectador no processo de aprendizagem e torna -se protagonista do próprio processo de ensino-aprendizagem.
Investigação Científica O eixo Investigação Científica que compõe o Itinerário na área de Ciência da Natureza tem como objetivo geral detalhar os conteúdos fundantes para subsidiar e interpretar os fenômenos, processos, dados com auxílio das abordagens investigativas, objetivando resolver situações cotidianas, com vista na promoção do desenvolvimento local e a melhoria da qualidade de vida ao qual o estudante está inserido.
Espera -se que, ao final deste eixo, os estudantes, por meio de ferramentas metodológicas que auxiliem na interlocução da interação, observação, levantamento de hipótese, realização de previsões e estimativas sejam capazes de empregar instrumentos de medição com base na interpretação de modelos, dados, resultados e processos que permitam avaliar e justificar suas conclusões no enfrentamento de situações-problema sob uma perspectiva científica.
Atrelados às práticas metodológicas, espera -se promover a participação social e o protagonismo dos estudantes como agentes de transformação de suas unidades de ensino e de suas comunidades, embasados em situação-problema que estejam no contexto do estudante, considerando demandas locais, regionais e globais, permitindo obter informações para compreender melhor o cenário, argumentar de forma crítica e apontar solução.
Processos Criativos O eixo Processos Criativos no Itinerário da Área de Ciências da Natureza e suas Tecnologias, propõe que os estudantes apliquem os conhecimentos as habilidades adquiridas de forma inovadora, usando sempre sua imaginação, propondo soluções que contribuam para resolução de \textbf{problemas} no Projeto de Vida e no mundo do trabalho.
Logo, deve fomentar a reflexão sobre elas e o processo criativo empregado, ou seja, assumam o papel de protagonismo de suas ações.
Por meio desse dinamismo de experiência criativa pretende -se que estudantes se tornem sujeitos autônomos, críticos, responsáveis.
Espera -se que, ao final, eles possam assumir os riscos das incertezas do mercado de trabalho e dos seus projetos de vida. \textbf{Mediação} e Intervenção Sociocultural O eixo \textbf{Mediação} e Intervenção Sociocultural enquanto parte integrante do Itinerário da área de Ciências da Natureza e suas Tecnologias propõe o envolvimento na vida pública via projetos de mobilização e intervenção sociocultural, por meio de atividades ancoradas em sequências didáticas investigativas que priorizem as interações dialógicas entre os envolvidos, associando os conceitos da CNT às Questões Sociocientíficas, Questões Controvérsias, explorando as abordagens da Ciência, Tecnologia, Sociedade e o Meio Ambiente ( CTSA ).
Logo, espera -se que as habilidades desenvolvidas ao longo deste eixo possam propiciar aos estudantes reconhecer e analisar questões sociais, culturais e \textbf{ambientais} diversas, considerando e incorporando a situação, a opinião, e o sentimento do outro, agindo com empatia flexibilidade e resiliência para promover o diálogo, a colaboração, a \textbf{mediação} e a resolução de conflitos, o combate ao preconceito e a valorização da diversidade. \textbf{Empreendedorismo} No eixo \textbf{Empreendedorismo} no Itinerário da Área de Ciências da Natureza e suas Tecnologias, sugere -se que os estudantes tenham foco estratégico nos seus empreendimentos e na elaboração de um bom planejamento nos seus projetos de vida e no mundo do trabalho.
Para tanto, busca desenvolver autonomia, foco e determinação para que consigam planejar e conquistar objetivos pessoais ou criar empreendimentos voltados à geração de renda via oferta de produtos e serviços, com ou sem uso de \textbf{tecnologias}.
Assim como, aprender a lidar de forma proativa e empreendedora com confiança para superar os grandes desafios e estarem preparados para superar situações de \textbf{estresses}, fracassos e \textbf{adversidades} na trajetória como cidadão social e consciente de seu papel na sociedade.
Os pressupostos metodológicos elencam as possíveis temáticas com suas respectivas abordagens didáticas, enquanto sugestão para serem exploradas nos eixos \textbf{estruturantes} dos Itinerários \textbf{Formativos}.
Os eixos \textbf{estruturantes} que integram os diferentes arranjos do Itinerário \textbf{Formativo} conectam experiências educativas do estudante com a realidade contemporânea, propiciando o desenvolvimento de habilidades para sua formação integral.
As sugestões para este Itinerário estão associadas a situações-problemas, por meio das abordagens apontadas, somadas ao uso de ferramentas pedagógicas – que devem basear -se na realidade local e escolar – espera -se que o Itinerário seja um mecanismo que propicie aos estudantes investigar, levantar hipóteses, dados, propor solução e comunicar a solução.
A temática escolhida deve, sobretudo, considerar o contexto em que os estudantes estão inseridos, e as especificidades do entorno da comunidade, podendo ser local, regional e/ou global.
Apresentamos no quadro abaixo, os pressupostos metodológicos com suas respectivas habilidades específicas, descritas por cada eixo \textbf{estruturante}, como possibilidades de temáticas para o Itinerário da área de Ciências da Natureza e suas Tecnologias.
É importante ressaltar que os pressupostos metodológicos permitem ao docente a criar diversas temáticas que podem ser estruturadas, dentro desse modelo, seguindo o contexto local, podendo cada eixo apresentar uma nova temática para trabalhar todos os eixos, ou um único assunto para trabalhar todos dos eixos no mesmo itinerário.

% matched lemmas: adversidade, ambiental, empreendedorismo, estresse, estruturante, formativo, mediação, problema, tecnologia
\end{textsample}
